\documentclass[10pt]{exam}

\usepackage[margin=1in]{geometry}
\usepackage{amsmath}
\usepackage{amssymb}
\usepackage{amsthm}
\usepackage{mathtools}
\usepackage{bm}
\usepackage{stmaryrd}

\usepackage{color}
\usepackage{colortbl}
\definecolor{deepblue}{rgb}{0,0,0.5}
\definecolor{deepred}{rgb}{0.6,0,0}
\definecolor{deepgreen}{rgb}{0,0.5,0}
\definecolor{gray}{rgb}{0.7,0.7,0.7}

\usepackage{hyperref}
\hypersetup{
  colorlinks   = true, %Colours links instead of ugly boxes
  urlcolor     = black, %Colour for external hyperlinks
  linkcolor    = blue, %Colour of internal links
  citecolor    = blue  %Colour of citations
}

\usepackage{listings}

%%%%%%%%%%%%%%%%%%%%%%%%%%%%%%%%%%%%%%%%%%%%%%%%%%%%%%%%%%%%%%%%%%%%%%%%%%%%%%%%

\theoremstyle{definition}
\newtheorem{problem}{Problem}
\newtheorem{example}{Example}
\newtheorem{defn}{Definition}
\newtheorem{fact}{Fact}
\newtheorem{refr}{References}
\newtheorem{theorem}{Theorem}
\newcommand{\E}{\mathbb E}
\newcommand{\R}{\mathbb R}
\DeclareMathOperator{\nnz}{nnz}
\DeclareMathOperator{\sign}{sign}
\DeclareMathOperator{\determinant}{det}
\DeclareMathOperator{\Var}{Var}
\DeclareMathOperator{\rank}{rank}
\DeclareMathOperator{\prob}{\mathbb P}
\DeclareMathOperator*{\argmin}{arg\,min}
\DeclareMathOperator*{\argmax}{arg\,max}

\newcommand{\Ein}{E_{\text{in}}}
\newcommand{\Eout}{E_{\text{out}}}
\newcommand{\Etest}{E_{\text{test}}}
\newcommand{\I}{\mathbf I}
\newcommand{\Q}{\mathbf Q}
\newcommand{\p}{\mathbf P}
\newcommand{\pb}{\bar {\p}}
\newcommand{\pbb}{\bar {\pb}}
\newcommand{\pr}{\bm \pi}

\newcommand{\trans}[1]{{#1}^{T}}
\newcommand{\loss}{\ell}
\newcommand{\w}{\mathbf w}
\newcommand{\wstar}{{\w}^{*}}
\newcommand{\x}{\mathbf x}
\newcommand{\y}{\mathbf y}
\newcommand{\lone}[1]{{\lVert {#1} \rVert}_1}
\newcommand{\ltwo}[1]{{\lVert {#1} \rVert}_2}
\newcommand{\lp}[1]{{\lVert {#1} \rVert}_p}
\newcommand{\linf}[1]{{\lVert {#1} \rVert}_\infty}
\newcommand{\lF}[1]{{\lVert {#1} \rVert}_F}

\newcommand{\mH}{m_{\mathcal H}}
\newcommand{\dvc}{{d_{\text{VC}}}}
\newcommand{\HH}[1]{\mathcal H_{\text{#1}}}
\newcommand{\Hbinary}{\HH_{\text{binary}}}
\newcommand{\Haxis}{\HH_{\text{axis}}}
\newcommand{\Hperceptron}{\HH_{\text{perceptron}}}


\newcommand{\ignore}[1]{}

%%%%%%%%%%%%%%%%%%%%%%%%%%%%%%%%%%%%%%%%%%%%%%%%%%%%%%%%%%%%%%%%%%%%%%%%%%%%%%%%

\begin{document}


\begin{center}
{
\Huge
%Chapter 1/2: The Learning Problem (III)
Chapter 2: Training vs Testing (II)
}
\end{center}

\begin{center}
%\includegraphics[height=3in]{scooby}
%~~~~~~~~~~
%\includegraphics[height=3in]{ml}
\end{center}


\begin{problem}
    What is the VC dimension of the following hypothesis classes?
    (These hypothesis classes all come from either Example 2.2 on page 43-45 or the problems in the back of the chapter.)

    \begin{enumerate}
    \item
        The set of positive rays in 1 dimension:
        \begin{equation}
            \HH{} = \bigg\{ x \mapsto \sign(x-a) : a \in \R \bigg\}
        \end{equation}
            \vspace{3in}

    \item
        The set of positive and negative rays in 1 dimension:
        \begin{equation}
            \HH{} = \bigg\{ x \mapsto \sigma\sign(x-a) : a \in \R, \sigma \in \{+1, -1\} \bigg\}
        \end{equation}
            \vspace{3in}

    \item The set of positive intervals in 1 dimension:
        \begin{equation}
            \HH{} = \bigg\{ x \mapsto \big\llbracket a \le x \le b \big\rrbracket : a \in \R, b \in \R \bigg\}
        \end{equation}
\vspace{4in}

    \item The set of positive and negative intervals in 1 dimension:
        \begin{equation}
            \HH{} = \bigg\{ x \mapsto \sigma \big\llbracket a \le x \le b \big\rrbracket : a \in \R, b \in \R, \sigma \in \{+1, -1\}  \bigg\}
        \end{equation}
\vspace{4in}

    \item The set of centered spheres in $\R^d$:
        \begin{equation}
            \HH{} = \bigg\{ \x \mapsto \big\llbracket \ltwo{\x} \le \alpha \big\rrbracket : \alpha \in \R^+ \bigg\}
        \end{equation}
\vspace{4in}

    %\item The set of non-centered spheres in $\R^d$:
        %\begin{equation}
            %\HH{} = \bigg\{ \x \mapsto \big\llbracket \ltwo{\x - \mu} \le \alpha \big\rrbracket : \alpha \in \R^+, \mu \in \R^d \bigg\}
        %\end{equation}
%\vspace{4in}
%
    %\item The set of non-centered spheres in $\R^d$:
        %\begin{equation}
            %\HH{} = \bigg\{ \x \mapsto \big\llbracket \ltwo{\x - \mu} \le \alpha \big\rrbracket : \alpha \in \R^+, \mu \in \R^d \bigg\}
        %\end{equation}
%\vspace{4in}

    \item Convex sets in 2 dimensions:
        $\mathcal H$ is the set of all functions $h : \R^2 \to \{+1, -1\}$ that are positive inside some convex set and negative elsewhere.
            (Recall that a set is convex if the line segment connecting any two points in the set lies entirely within the set.)
\vspace{4in}

    %\item The set of decision stumps:
    %\begin{equation}
    %\end{equation}
    \end{enumerate}
\end{problem}

%\begin{problem}
%Recall that the textbook uses $f$ to denote the \emph{true label function}, which can defined to be
%\begin{equation}
    %f = \argmin_{h \in \mathcal H^*} \Eout(h),
%\end{equation}
%where $\mathcal H^*$ is the union of all hypothesis classes.
%\begin{enumerate}
    %\item According to VC-theory, why does it not make sense to use $H^*$ as our hypothesis class?
%\end{enumerate}
%\end{problem}

%\begin{problem}
%\end{problem}

%\begin{defn}
    %The \emph{Bayes Optimal Predictor} $f$
    %\begin{equation}
        %f = \argmin_h \Eout(h)
    %\end{equation}
%\end{defn}

%\newpage
\begin{problem}
    For each statement below, indicate whether it is true or false and explain why.
    (The best possible explanation for true answers is a proof, and the best possible explanation for false answers is a counterexample.)
    \begin{enumerate}
        \item Let $f$ be the true labeling function.  Then for all data distributions, $\Eout(f) = 0$.
            \vspace{3in}

        \item Let $\mathcal H$ be the perceptron hypothesis class, and $g$ be the result of the PLA.
            Also let $f$ be the true labeling function.
            Then $\Eout(g) \ge \Eout(f)$.
            \vspace{3in}

            \newpage
        \item Let $\mathcal H$ be the perceptron hypothesis class, and $g$ be the result of the PLA.
            Also let $f$ be the true labeling function.
            Then $\Ein(g) \ge \Ein(f)$.
            \vspace{3in}

        \item Let $f$ be the true labeling function.
            There exists some hypothesis class $\mathcal H$
            with hypothesis $g \in \mathcal H$ satisfying
            $\Eout(g) < \Eout(f)$.
            \vspace{3in}

        \item If $\mathcal H$ is a finite hypothesis class, and $g$ is trained using the TEA algorithm, the it is always true that $|\Ein(g) - \Eout(g)| \le \Ein(g)$.
            \vspace{3in}

        \item Let $g$ be a hypothesis selected from some hypothesis class $\mathcal H$ with finite dimension.  Then it is possible for $\Ein(g)$ to be less than $\Eout(g)$.
            \vspace{3in}

        \item The VC dimension of every finite hypothesis class is finite.
            \vspace{3in}
        \item The VC dimension of every infinite hypothesis class is infinite.
            \vspace{3in}
        %\item There exists a hypothesis $h$ that has lower error than the Bayes Optimal Predictor.
        \item If there exists a hypothesis $h\in\mathcal H$ such that $\min_{h\in\mathcal H} \Ein(h) = 0$,
            then the VC dimension of $\mathcal H$ must be finite.
            \vspace{3in}
        %\item Let $\mathcal H$ be a finite hypothesis class of size $M$.
            %Then $\mH = O(M)$.
        \item
            If $\mathcal H$ can shatter a set of size $N$, then $\dvc(\mathcal H) \ge N$.
            \vspace{3in}
        \item
            For any hypothesis class $\mathcal H$, the value $\dvc(\mathcal H) + 1$ is a break point for $\mathcal H$.
            \vspace{3in}
        \item
            If the hypothesis class $\mathcal H$ shatters some dataset of size $N$, then $\mH(N) = 2^N$.
            \vspace{3in}
        \item
            There exists some hypothesis class $\mathcal H$ with growth function $\mH(N) = \Theta(2^{\sqrt{N}})$.
            \vspace{3in}
        \item
            Let $\mathcal H$ be an arbitrary hypothesis class.
            Then $\mH(N) = O(2^N)$.
            \vspace{3in}
        \item
            Let $\mathcal H_1$ and $\mathcal H_2$ be hypothesis classes satisfying $\mathcal H_1 \subset \mathcal H_2$.
            Then $\dvc(\mathcal H_1) \le \dvc(\mathcal H_2)$.
            \vspace{3in}
        \item
            Let $\mathcal H_1$ and $\mathcal H_2$ be hypothesis classes satisfying $\dvc(\mathcal H_1) \le \dvc(\mathcal H_2)$.
            Then $\mathcal H_1 \subset \mathcal H_2$.
            \vspace{3in}
        \item
            Let $\mathcal H_1$ and $\mathcal H_2$ be hypothesis classes satisfying $\dvc(\mathcal H_1) = \dvc(\mathcal H_2)$.
            Then $\mathcal H_1 = \mathcal H_2$.
            \vspace{3in}
        %\item
            %If the input data is $d$ dimensional, then the VC dimension of every hypothesis class must be at least $d$.
            %\vspace{3in}
        %%\item
            %%If the hypothesis class $\mathcal H$ shatters some dataset of size $N$, then $\mH(N) = 2^N$.
        %\item Finite hypothesis classes of size $M$ can have VC dimension $O(1)$.
            %\vspace{3in}
    \end{enumerate}
\end{problem}

%\begin{problem}
    %You
%\end{problem}

\end{document}



